\documentclass[10pt,a4paper]{article}
\usepackage[utf8]{inputenc}
\usepackage[T1]{fontenc}
\usepackage{lmodern}
\usepackage[margin=1in]{geometry}
\usepackage{xcolor}
\usepackage{amsmath,amssymb,amsfonts}
\usepackage{graphicx}
\usepackage{booktabs}
\usepackage{adjustbox}
\usepackage{tcolorbox}
\usepackage{enumitem}
\usepackage{microtype}
\usepackage{array}
\usepackage{fancyhdr}

% Custom Colors matching house style
\definecolor{bitsnavy}{RGB}{33,53,94}
\definecolor{bitsblue}{RGB}{44,90,160}
\definecolor{bitsgreen}{RGB}{46,125,82}
\definecolor{bitsamber}{RGB}{200,136,30}
\definecolor{bitsred}{RGB}{178,58,72}

% Callout Boxes
\tcolorboxenvironment{intuition}{colback=bitsblue!8,colframe=bitsblue,title=\textbf{Everyday Picture \& Intuition},fonttitle=\bfseries,boxrule=0.8pt,arc=3pt}
\tcolorboxenvironment{worked}{colback=bitsgreen!8,colframe=bitsgreen,title=\textbf{Fully Worked Example},fonttitle=\bfseries,boxrule=0.8pt,arc=3pt}
\tcolorboxenvironment{everyday}{colback=bitsnavy!5,colframe=bitsnavy,title=\textbf{Real-World Picture},fonttitle=\bfseries,boxrule=0.8pt,arc=3pt}
\tcolorboxenvironment{watchout}{colback=bitsred!8,colframe=bitsred,title=\textbf{Watch Out! Pitfalls},fonttitle=\bfseries,boxrule=0.8pt,arc=3pt}
\tcolorboxenvironment{keytake}{colback=bitsamber!10,colframe=bitsamber,title=\textbf{Key Takeaways},fonttitle=\bfseries,boxrule=0.8pt,arc=3pt}

\newenvironment{intuition}{\begin{tcolorbox}}{\end{tcolorbox}}
\newenvironment{worked}{\begin{tcolorbox}}{\end{tcolorbox}}
\newenvironment{everyday}{\begin{tcolorbox}}{\end{tcolorbox}}
\newenvironment{watchout}{\begin{tcolorbox}}{\end{tcolorbox}}
\newenvironment{keytake}{\begin{tcolorbox}}{\end{tcolorbox}}

\newenvironment{steps}{\begin{enumerate}[label=\textbf{Step \arabic*:},leftmargin=*]}{\end{enumerate}}

\newcommand{\secsub}[1]{\noindent\textit{\large #1}\vspace{0.8em}}
\newcommand{\housefig}[2]{
\begin{figure}[htbp]
\centering
\includegraphics[width=0.85\linewidth]{#1}
\caption{#2}
\end{figure}
}

\pagestyle{fancy}
\fancyhf{}
\rhead{\color{bitsnavy}\bfseries WILP, BITS Pilani}
\lhead{\color{bitsnavy}\bfseries ISM Session 13 Study Companion}
\rfoot{Page \thepage}

\begin{document}

% Banner
\begin{tcolorbox}[colback=bitsnavy,colframe=bitsnavy,arc=0pt,outer arc=0pt,top=12pt,bottom=12pt]
\color{white}
\centering
{\large\bfseries WILP, BITS Pilani}\\[4pt]
{\LARGE\bfseries Session 13: Time Series Forecasting Methods}\\[4pt]
{\small Course: Introduction to Statistical Methods (ISM) \quad|\quad Instructor: Dr. YVK Ravi Kumar}
\end{tcolorbox}

\vspace{1em}

\section{Moving Averages \& Forecast Error Accuracy}
\secsub{Smooth out short-term fluctuations and evaluate forecast accuracy using MSE, MAD, MAPE.}

Moving averages assume that future observations will resemble recent historical levels. Simple Moving Averages (SMA) weigh all $k$ recent periods equally, whereas Weighted Moving Averages (WMA) assign larger weights to more recent observations. Forecast error is $e_t = y_t - F_t$.

\begin{intuition}
A simple moving average is like taking average temperature over the last 3 days to guess tomorrow's weather. Accuracy metrics like MSE (Mean Squared Error) and MAD (Mean Absolute Deviation) measure how far off your guesses are on average.
\end{intuition}

\subsection{Mathematical Formulas}
\begin{align}
\text{SMA}_t &= \frac{y_{t-1} + y_{t-2} + \dots + y_{t-k}}{k}, \quad \text{WMA}_t = \frac{w_1 y_{t-1} + w_2 y_{t-2} + \dots + w_k y_{t-k}}{\sum w_i} \\
\text{MSE} &= \frac{1}{n} \sum (y_t - F_t)^2, \quad \text{MAD} = \frac{1}{n} \sum |y_t - F_t|, \quad \text{MAPE} = \frac{100\%}{n} \sum \left|\frac{y_t - F_t}{y_t}\right|
\end{align}

\housefig{figures/moving_average.png}{Comparison of raw production data against a 3-period moving average curve.}

\begin{worked}
Demand observations for recent months: Month 49 = 130, Month 50 = 128, Month 51 = 137.
\begin{steps}
\item Compute 3-month Simple Moving Average forecast for Month 52:
\begin{equation}
F_{52} = \frac{137 + 128 + 130}{3} = \frac{395}{3} = 131.67
\end{equation}
\item Compute Weighted Moving Average with weights $w = [3, 2, 1]$ (highest weight on most recent period):
\begin{equation}
F_{52} = \frac{3(137) + 2(128) + 1(130)}{3 + 2 + 1} = \frac{411 + 256 + 130}{6} = \frac{797}{6} = 132.83
\end{equation}
\end{steps}
\end{worked}

\section{Simple Exponential Smoothing}
\secsub{Assign exponentially decaying weights to past historical observations.}

Simple Exponential Smoothing (SES) updates forecasts using a smoothing constant $\alpha \in [0, 1]$:
\begin{equation}
F_{t+1} = \alpha y_t + (1 - \alpha) F_t = F_t + \alpha (y_t - F_t)
\end{equation}

\begin{intuition}
Exponential smoothing is like updating your guess for commute time. Your new guess equals your old guess plus a fraction $\alpha$ of your error from today's commute.
\end{intuition}

\housefig{figures/exponential_smoothing.png}{Sensitivity of exponential smoothing forecasts under different smoothing constants $\alpha = 0.2$ vs $\alpha = 0.8$.}

\begin{worked}
Suppose forecast for the first week of March was $F_1 = 500$ units, but actual demand was $y_1 = 450$ units. Let $\alpha = 0.3$.
\begin{steps}
\item Calculate forecast for week 2:
\begin{equation}
F_2 = 0.3(450) + (1 - 0.3)(500) = 135 + 350 = 485
\end{equation}
\item Suppose actual demand in week 2 is $y_2 = 505$ units. Calculate forecast for week 3:
\begin{equation}
F_3 = 0.3(505) + 0.7(485) = 151.5 + 339.5 = 491.0
\end{equation}
\end{steps}
\end{worked}

\section{Holt's Double Exponential Smoothing}
\secsub{Track both baseline level and linear trend simultaneously in trending series.}

When data exhibits a persistent linear trend, simple exponential smoothing lags behind. Holt's Double Exponential Smoothing introduces a second smoothing equation to update the trend component $T_t$ alongside level $L_t$.

\begin{intuition}
If sales grow by 10 units every month, simple smoothing will always underestimate next month's sales. Holt's method tracks both current height and growth speed.
\end{intuition}

\subsection{Holt's Equations}
\begin{align}
L_t &= \alpha y_t + (1 - \alpha)(L_{t-1} + T_{t-1}) \\
T_t &= \beta (L_t - L_{t-1}) + (1 - \beta) T_{t-1} \\
\hat{y}_{t+h} &= L_t + h T_t
\end{align}

\housefig{figures/holt_double_exp.png}{Holt's double exponential smoothing tracking trending time series without lagging.}

\begin{watchout}
Simple moving averages and single exponential smoothing are only suitable for stationary series. For trending data, use Holt's method or ARIMA.
\end{watchout}

\begin{keytake}
\begin{itemize}
\item Simple Moving Averages average recent $k$ values equally; WMAs emphasize recent data.
\item Error metrics MSE, MAD, and MAPE measure forecast accuracy.
\item Exponential Smoothing updates forecasts by adding a fraction $\alpha$ of past forecast error.
\item Holt's Double Exponential Smoothing incorporates a trend factor to eliminate lag on trending data.
\end{itemize}
\end{keytake}

\end{document}
